\PassOptionsToPackage{dvipsnames}{xcolor}
\documentclass[11pt,letterpaper]{article}
\usepackage{mathalpha}
% Packages specific to this document
\usepackage{algorithm}
\usepackage{algorithmicx}
\usepackage{algpseudocode}
\usepackage{graphicx}
\usepackage{subcaption}
\usepackage{booktabs}
\usepackage{float}
\usepackage{bm}
\usepackage{mathtools}
\usepackage{multirow} 
\usepackage{mathrsfs}
\usepackage{nicefrac}
\usepackage{listings}
\usepackage{dsfont}
\usepackage{fdsymbol}
\usepackage{wrapfig}
\usepackage{array}
\usepackage{multicol}
\usepackage{url}
\usepackage{fancyhdr}
\usepackage{rotating}
\usepackage[capitalise]{cleveref}
\usepackage{tikz}
\usetikzlibrary{shapes,arrows,positioning,calc}

% Citation and bibliography setup
\usepackage[comma,numbers,compress]{natbib}
\bibliographystyle{plainnat}

% Caption setup
\usepackage[font=small,labelfont=bf]{caption}

% Color definitions
\usepackage{xcolor}
\definecolor{lightblue}{rgb}{0.22,0.45,0.70}
\definecolor{YaleBlue}{rgb}{0.06,0.30,0.56}
\definecolor{memstblue}{RGB}{41, 98, 255}

% Hyperref setup - load hypcap after hyperref
\usepackage{hyperref}
\hypersetup{
    colorlinks = true,
    citecolor = {YaleBlue},
}
\usepackage{hypcap}

\title{\centering \textbf{MemSt}: A Git-Like Semantic Memory Architecture for AI Agents}
\author{\centering Yifan Yang\\
\vspace{1em}
\hspace{0.6cm} \href{mailto:yfyang.86@hotmail.com}{\texttt{yfyang.86@hotmail.com}}}

% Setup page style
\pagestyle{fancy}
\fancyhf{}
\fancyhead[C]{\textit{\footnotesize MemSt: Git-Like Semantic Memory for AI Agents}}
\fancyfoot[C]{\thepage}

\begin{document}

\maketitle

% Abstract
\begin{abstract}
\vspace{1cm}

Large Language Models (LLMs) have demonstrated remarkable capabilities in generating contextually coherent responses, yet their fixed context windows and ephemeral memory pose fundamental challenges for maintaining consistency over prolonged multi-session interactions. We introduce \textbf{MemSt}, a novel semantic memory architecture that addresses these limitations through a Git-inspired content-addressable storage system with tiered memory management and knowledge graph extraction.

MemSt provides a unique combination of features not found in existing memory systems: (1) \textit{Git-like versioning} with branch, merge, and blame operations on memory states; (2) \textit{Four-tier memory hierarchy} (Working, Short-term, Long-term, Archival) with automatic lifecycle management; (3) \textit{Hybrid search} combining BM25, HNSW vector search, and Reciprocal Rank Fusion; (4) \textit{Knowledge Graph Extraction v2} with 80+ domain ontologies and multi-backend storage (SQLite/DuckDB); (5) \textit{Sleep-time consolidation} for asynchronous memory maintenance; and (6) \textit{Multi-agent worktrees} for isolated concurrent agent execution.

Our architecture supports multiple interfaces including native Rust library, Python bindings (PyO3), REST API (FastAPI), and React-based Web UI with integrated Nanobot agent framework. Through comprehensive architectural evaluation, we demonstrate that MemSt achieves superior memory organization, faster retrieval latency, and more efficient token utilization compared to existing memory-augmented systems.

\vspace{5mm}
\textbf{Code}: \href{https://github.com/yfyang86/memst}{https://github.com/yfyang86/memst}

\end{abstract}

% Sections
\section{Introduction}
\label{sec:intro}

Human memory is a \emph{foundation of intelligence}---it shapes our identity, guides decision-making, and enables us to learn, adapt, and form meaningful relationships \citep{craik1992human}. Among its many roles, memory is essential for communication: we recall past interactions, infer preferences, and construct evolving mental models of those we engage with \citep{assmann2011communicative}. This ability to retain and retrieve information over extended periods enables coherent, contextually rich exchanges that span days, weeks, or even months.

AI agents, powered by large language models (LLMs), have made remarkable progress in generating fluent, contextually appropriate responses \citep{yu2024finmem, zhang2024survey}. However, these systems are fundamentally limited by their reliance on fixed context windows, which severely restrict their ability to maintain coherence over extended interactions \citep{bulatov2022recurrent, liu2023think}. This limitation stems from LLMs' lack of persistent memory mechanisms that can extend beyond their finite context windows.

The absence of persistent memory creates a fundamental disconnect in human-AI interaction. Without memory, AI agents forget user preferences, repeat questions, and contradict previously established facts. Consider a scenario where a user mentions being vegetarian and avoiding dairy products in an initial conversation. In a subsequent session, when the user asks about dinner recommendations, a system without persistent memory might suggest chicken, completely contradicting the established dietary preferences.

Beyond conversational settings, memory mechanisms have been shown to dramatically enhance agent performance in interactive environments \citep{majumderclin, shinn2023reflexion}. Agents equipped with memory of past experiences can better anticipate user needs, learn from previous mistakes, and generalize knowledge across tasks \citep{chhikara2023knowledge}. Research demonstrates that memory-augmented agents improve decision-making by leveraging causal relationships between actions and outcomes.

\subsection{Current Limitations of Memory Systems}

Existing memory solutions for LLM agents fall into several categories, each with significant limitations:

\textbf{Vector databases} (Chroma, Milvus, Pinecone) store embeddings of conversation history but lack structured organization, versioning, and lifecycle management. They treat memory as an undifferentiated ``bag of vectors,'' making it difficult to maintain coherent narrative structures or track memory evolution over time.

\textbf{Cloud memory services} (Mem0, Letta) provide managed memory but require external dependencies, introducing latency, cost, and data sovereignty concerns. They typically offer limited or no versioning capabilities, making it impossible to audit memory changes or revert to previous states.

\textbf{In-context memory} approaches simply prepend conversation history to the prompt. While simple, they suffer from quadratic attention costs and inevitable context overflow as conversations extend.

\textbf{Graph memory systems} (Zep, A-MEM) capture entity relationships but often lack efficient retrieval mechanisms, multi-tier organization, or local deployment options. They may also suffer from high token overhead and slow construction times.

\subsection{Key Challenges}

We identify four critical challenges that existing memory systems fail to address:

\begin{enumerate}
    \item \textbf{No Versioning}: Memory is mutable state without history. Agents cannot revert to previous memory states, branch experiments, or audit when specific facts were learned.
    
    \item \textbf{No Lifecycle Management}: Memories lack automatic promotion, expiration, or consolidation. Working memories are never summarized; stale memories are never archived.
    
    \item \textbf{No Token-Budget Awareness}: Retrieval does not respect LLM context limits. Systems fetch memories without considering whether they fit within available token budgets.
    
    \item \textbf{No Multi-Agent Isolation}: Concurrent agents share memory space, risking corruption and cross-contamination of knowledge.
\end{enumerate}

\subsection{Our Contribution: MemSt}

We introduce \textbf{MemSt} (pronounced ``mem-st''), a semantic memory architecture inspired by Git's content-addressable storage model. MemSt treats memory as a versioned, structured repository rather than an ephemeral cache. Our evaluation on the LOCOMO benchmark demonstrates that MemSt achieves \textbf{49.2\% F1} with the QuestionAdaptive strategy---a \textbf{118\% improvement} over standard RAG (22.6\%) and outperforming published systems including Mem0 (34.8\%) and Zep (35.7\%).

Our key contributions include:

\begin{itemize}
    \item \textbf{Git-Like Memory Repository}: Content-addressable storage with Blake3 hashing, branch and merge operations, and full provenance tracking via a write-ahead log (WAL).
    
    \item \textbf{Hierarchical Retrieval Strategies}: Multi-level memory organization (session summaries $\rightarrow$ key turns $\rightarrow$ full context) with dynamic token budget allocation, achieving 41.7\% F1 with 85ms latency.
    
    \item \textbf{Specialized Graph Architectures}: TemporalKG for time-aware queries (54.0\% F1, +120\% vs RAG) and MultiHopKG for relationship chains (52.8\% F1, +133\% vs RAG).
    
    \item \textbf{Cascade Retrieval}: Three-stage filtering (BM25 $\rightarrow$ Vector $\rightarrow$ Cross-encoder) achieving 55.2\% F1 on single-hop questions.
    
    \item \textbf{Question-Adaptive Routing}: Dynamic strategy selection based on question classification, achieving best overall performance (49.2\% F1) by combining specialized strategies.
    
    \item \textbf{Optimized Hybrid Search}: Category-tuned Reciprocal Rank Fusion with BM25 parameters ($k_1=0.5$, $b=0.75$) optimized for conversational text.
    
    \item \textbf{Four-Tier Memory Hierarchy}: Working (hot), Short-term (recent), Long-term (important), and Archival (historical) tiers with automatic lifecycle management.
    
    \item \textbf{Production-Ready Implementation}: Pure Rust core with 3.2$\times$ lower latency than RAG, Python bindings, REST API, and MCP adapter for seamless integration.
\end{itemize}

MemSt is designed for production deployment: pure Rust core for performance, local-first architecture for data sovereignty, and comprehensive APIs for integration with existing agent frameworks.

\section{System Architecture}
\label{sec:architecture}

MemSt adopts a layered architecture (Figure \ref{fig:architecture}) with clear separation between interfaces, language bindings, and core storage components. This section describes the key architectural components.

\begin{figure*}[t]
    \centering
    \begin{tikzpicture}[
        node distance=1.5cm,
        box/.style={rectangle, draw, rounded corners, minimum width=3cm, minimum height=0.8cm, align=center, fill=lightblue!20},
        arrow/.style={->, >=stealth, thick}
    ]
        % Top layer
        \node[box, fill=memstblue!30] (ui) {Web UI / React};
        \node[box, fill=memstblue!30, right=of ui] (api) {REST API / FastAPI};
        \node[box, fill=memstblue!30, right=of api] (mcp) {MCP Adapter};
        \node[box, fill=memstblue!30, right=of mcp] (cli) {CLI / memst};
        
        % Middle layer
        \node[box, below=of ui, xshift=1.5cm] (python) {Python Bindings (PyO3)};
        \node[box, right=of python, xshift=1.5cm] (rust) {Rust Core (memst-core)};
        
        % Bottom layers
        \node[box, below=of python, fill=green!20] (search) {Hybrid Search\\BM25 + HNSW};
        \node[box, right=of search, fill=green!20] (repo) {MemRepo\\Git-like Store};
        \node[box, right=of repo, fill=green!20] (kg) {KG Extraction v2\\SQLite/DuckDB};
        \node[box, right=of kg, fill=green!20] (sleep) {Sleep-Time\\Consolidation};
        
        % Arrows
        \draw[arrow] (ui) -- (python);
        \draw[arrow] (api) -- (python);
        \draw[arrow] (mcp) -- (python);
        \draw[arrow] (cli) -- (rust);
        \draw[arrow] (python) -- (rust);
        \draw[arrow] (rust) -- (search);
        \draw[arrow] (rust) -- (repo);
        \draw[arrow] (rust) -- (kg);
        \draw[arrow] (rust) -- (sleep);
    \end{tikzpicture}
    \caption{MemSt architecture overview showing the multi-layer design with interfaces (top), language bindings (middle), and core storage/processing components (bottom).}
    \label{fig:architecture}
\end{figure*}

\subsection{Git-Like Memory Repository (MemRepo)}

At the foundation of MemSt is \textbf{MemRepo}, a content-addressable storage system inspired by Git's object model. Unlike traditional databases that use location-based addressing, MemRepo stores all data as immutable objects referenced by their Blake3 hash.

\subsubsection{Object Model}

MemRepo implements four fundamental object types:

\begin{itemize}
    \item \textbf{Blob}: Stores arbitrary binary content (messages, memories, embeddings). Each blob is addressed by $\text{Blake3}(\text{content})$.
    
    \item \textbf{Tree}: Represents hierarchical directory structures, mapping names to blob hashes and subtree hashes. Enables organized storage of session data.
    
    \item \textbf{Commit}: Captures a snapshot of the memory state at a point in time, containing a tree hash, parent commit hash(es), author/timestamp metadata, and a commit message.
    
    \item \textbf{Tag}: Named references to specific commits, enabling semantic labeling of memory milestones (e.g., ``session-start'', ``checkpoint-before-merge'').
\end{itemize}

This object model provides several advantages: (1) \textit{deduplication}---identical content stored once regardless of location; (2) \textit{immutability}---historical states are permanently preserved; (3) \textit{integrity}---any corruption is immediately detectable via hash mismatch.

\subsubsection{Write-Ahead Log (WAL)}

For durability and crash recovery, MemRepo implements a write-ahead log:

\begin{enumerate}
    \item All mutations are first appended to the WAL as journal entries
    \item Entries include operation type, affected objects, and rollback information
    \item Periodic checkpoints flush committed state to the object store
    \item On recovery, replay WAL entries from last checkpoint
\end{enumerate}

The WAL ensures ACID semantics: operations are atomic (all-or-nothing), consistent (integrity constraints enforced), isolated (via versioning), and durable (WAL persistence).

\subsubsection{Branching and Merging}

MemRepo supports Git-style branching for memory evolution:

\begin{itemize}
    \item \textbf{Branch}: A named reference to a commit, allowing parallel memory evolution
    \item \textbf{Checkout}: Switch the active branch to a different memory state
    \item \textbf{Merge}: Combine changes from two branches with three-way merge semantics
\end{itemize}

Branching enables powerful workflows: experimental memory configurations can be tested in isolation, agent behaviors can be versioned and rolled back, and multi-agent scenarios can use per-agent branches.

\subsection{Four-Tier Memory Hierarchy}

MemSt organizes memories into four tiers based on access frequency and importance:

\begin{table}[h]
\centering
\caption{Memory Tier Characteristics}
\begin{tabular}{@{}lllll@{}}
\toprule
\textbf{Tier} & \textbf{Access} & \textbf{Capacity} & \textbf{Promotion} & \textbf{Eviction} \\
\midrule
Working & Fastest & Limited & Manual & Age/importance \\
Short-term & Fast & Moderate & Access count & TTL \\
Long-term & Slower & Large & Importance score & Archival \\
Archival & Slowest & Unlimited & N/A & Manual \\
\bottomrule
\end{tabular}
\label{tab:memory-tiers}
\end{table}

\subsubsection{Working Memory}

Working memory contains the active conversation context and pinned facts. It is the only tier directly injected into the LLM context window. Working memory has strict token budget enforcement:

\begin{equation}
\sum_{m \in \text{Working}} \text{tokens}(m) \leq B_{\text{working}}
\end{equation}

where $B_{\text{working}}$ is configurable (default: 2048 tokens). When the budget is exceeded, memories are promoted to Short-term based on recency and importance scores.

\subsubsection{Short-Term Memory}

Short-term memory contains recent facts and user preferences that may be needed soon but are not immediately relevant. It is stored in compressed binary format with BM25 indexing for fast retrieval.

Promotion to Short-term occurs when:
\begin{itemize}
    \item Working memory exceeds its token budget
    \item A memory has been accessed more than $A_{\min}$ times (default: 3)
\end{itemize}

Demotion to Long-term occurs when a memory's time-to-live (TTL) expires (default: 7 days).

\subsubsection{Long-Term Memory}

Long-term memory stores important knowledge and historical data. Memories reach Long-term based on:

\begin{equation}
I(m) = \alpha \cdot \text{access\_freq}(m) + \beta \cdot \text{user\_rating}(m) + \gamma \cdot \text{semantic\_density}(m)
\end{equation}

where $I(m)$ is the importance score and $\alpha, \beta, \gamma$ are configurable weights. Memories with $I(m) > \theta_{\text{long}}$ are promoted to Long-term.

\subsubsection{Archival Memory}

Archival memory contains rarely accessed historical data. It is stored in highly compressed format and may be offloaded to slower storage (e.g., network-attached storage, S3). Archival memories retain full semantic indexing for retrieval but with higher latency.

\subsection{Context Assembly with Token Budget}

A unique feature of MemSt is \textbf{token-budget-aware context assembly}. When building the LLM prompt, MemSt does not simply retrieve the top-$k$ memories; it solves a constrained optimization problem:

\begin{align}
\text{maximize} \quad & \sum_{m \in S} R(m, q) \cdot I(m) \\
\text{subject to} \quad & \sum_{m \in S} \text{tokens}(m) \leq B_{\text{context}} - B_{\text{reserved}} \\
& |S| \leq N_{\max}
\end{align}

where $R(m, q)$ is the relevance of memory $m$ to query $q$, $I(m)$ is importance, $B_{\text{context}}$ is the model's context window, $B_{\text{reserved}}$ is tokens reserved for the response, and $N_{\max}$ is a maximum memory count for diversity.

This optimization ensures that the context window is filled with the most valuable information possible, rather than truncating an oversized retrieval set.

\subsection{Hybrid Search}

MemSt implements a three-backend search system:

\subsubsection{Native BM25}

A pure Rust implementation of the BM25 ranking function:

\begin{equation}
\text{BM25}(q, d) = \sum_{t \in q} \text{IDF}(t) \cdot \frac{f(t, d) \cdot (k_1 + 1)}{f(t, d) + k_1 \cdot (1 - b + b \cdot \frac{|d|}{\text{avgdl}})}
\end{equation}

with default parameters $k_1 = 1.2$, $b = 0.75$. This backend requires no external dependencies and is suitable for edge deployment.

\subsubsection{HNSW Vector Search}

For semantic similarity, MemSt uses Hierarchical Navigable Small World (HNSW) graphs:

\begin{itemize}
    \item Embeddings are computed via configurable API (OpenAI, local models)
    \item HNSW index enables approximate nearest neighbor search in $O(\log N)$
    \item Supports cosine similarity, Euclidean distance, and dot product
\end{itemize}

\subsubsection{Reciprocal Rank Fusion (RRF)}

To combine keyword and semantic signals, MemSt uses RRF:

\begin{equation}
\text{RRF}(d) = \sum_{r \in R} \frac{1}{k + r_r(d)}
\end{equation}

where $R$ is the set of rankers (BM25, HNSW), $r_r(d)$ is document $d$'s rank in ranker $r$, and $k = 60$ is a constant to dampen the impact of high rankings. This fusion approach is robust and requires no training.

\subsection{Multi-Agent Worktrees}

Inspired by Git worktrees, MemSt provides memory isolation for concurrent agents:

\begin{itemize}
    \item Each agent gets its own \textbf{worktree}---a separate working directory backed by the same MemRepo
    \item Worktrees can be created, switched, and removed without affecting other agents
    \item Changes in one worktree are invisible to others until explicitly merged
    \item The main worktree serves as a shared baseline
\end{itemize}

This architecture enables safe multi-agent deployment: agent A can experiment with memory modifications while agents B and C operate on stable baselines, with explicit merge gates controlling integration.

\section{Knowledge Graph Extraction v2}
\label{sec:kg_extraction}

Beyond flat memory storage, MemSt incorporates \textbf{Knowledge Graph Extraction v2} (KGv2), a structured knowledge representation system that extracts entities and relationships from conversations using LLM-powered extraction.

\subsection{Design Motivation}

While tiered memory excels at storing and retrieving factual statements, it lacks explicit relational structure. Consider the following conversation:

\begin{quote}
\textit{``Alice moved to San Francisco last year. She works at OpenAI as a research scientist. Her manager is Bob, who previously worked at Google.''}
\end{quote}

A flat memory system might store three separate facts:
\begin{itemize}
    \item Alice moved to San Francisco in 2022
    \item Alice works at OpenAI
    \item Bob is Alice's manager
\end{itemize}

But the \textit{relationships} between these facts---that Bob works at OpenAI, that Alice's move and job change are related, that Bob's Google experience influences his management---are implicit at best.

KGv2 makes these relationships explicit, enabling:
\begin{enumerate}
    \item \textbf{Multi-hop reasoning}: Answering ``Where did Alice's manager work before?'' requires traversing Alice $\rightarrow$ manager $\rightarrow$ Bob $\rightarrow$ previously worked at $\rightarrow$ Google.
    \item \textbf{Temporal reasoning}: Tracking when relationships were established and how they evolved.
    \item \textbf{Conflict detection}: Identifying when new information contradicts existing relationships.
\end{enumerate}

\subsection{Architecture Overview}

KGv2 follows a modular pipeline architecture:

\begin{enumerate}
    \item \textbf{Extraction}: LLM extracts entities and relationship triplets from text
    \item \textbf{Linking}: Mentions of the same entity are resolved to canonical nodes
    \item \textbf{Storage}: Entities and relationships are persisted to the backend
    \item \textbf{Retrieval}: Query-time subgraph extraction for context assembly
\end{enumerate}

\subsection{Multi-Backend Storage}

A key innovation of KGv2 is pluggable storage backends:

\subsubsection{SQLite Backend}

The default backend uses SQLite with the schema:

\begin{verbatim}
entities(id, name, entity_type, embedding, created_at)
relationships(id, source_id, target_id, relation_type, 
              confidence, created_at, valid_until)
ontologies(id, top_category, first_category, 
           second_category, description)
\end{verbatim}

SQLite is lightweight, requires no external services, and is suitable for development and edge deployment.

\subsubsection{DuckDB Backend}

For production analytics workloads, KGv2 supports DuckDB:

\begin{itemize}
    \item Vectorized query execution for analytical workloads
    \item Efficient aggregations and graph analytics (PageRank, community detection)
    \item Parquet export for integration with data lakes
\end{itemize}

The backend is selected at compile time via Cargo features:

\begin{verbatim}
# SQLite (default)
memst-extract-v2 = { path = "../memst-extract-v2" }

# DuckDB (production)
memst-extract-v2 = { path = "../memst-extract-v2", 
                     default-features = false, 
                     features = ["duckdb"] }
\end{verbatim}

\subsection{Ontology System}

KGv2 includes a comprehensive ontology system supporting 80+ intelligence domains:

\begin{table}[ht]
\centering
\caption{Sample Ontology Categories}
\begin{tabular}{@{}lll@{}}
\toprule
\textbf{Top Category} & \textbf{First Category} & \textbf{Second Category} \\
\midrule
Domain Intelligence & Tech Intelligence & Artificial Intelligence \\
Domain Intelligence & Tech Intelligence & Semiconductor \\
Domain Intelligence & Business Intelligence & Market Analysis \\
Domain Intelligence & Business Intelligence & Competitive Intelligence \\
Social Intelligence & Social Media & Trend Analysis \\
Social Intelligence & Public Opinion & Sentiment Monitoring \\
Geopolitical & Regional & Asia-Pacific \\
Geopolitical & Regional & European Union \\
\bottomrule
\end{tabular}
\label{tab:ontologies}
\end{table}

Each ontology defines:
\begin{itemize}
    \item Expected entity types (Person, Organization, Location, Event, Technology)
    \item Relationship vocabulary (works\_at, located\_in, acquired, partnered\_with)
    \item Extraction prompts tuned for the domain
\end{itemize}

Users can load custom ontologies via JSON schema:

\begin{verbatim}
{
  "top_category": "Custom Domain",
  "first_category": "Research",
  "second_category": "Machine Learning",
  "chinese_name": "Research-ML",
  "english_name": "Research-ML",
  "overview": "ML research tracking"
}
\end{verbatim}

\subsection{Entity Extraction Pipeline}

The extraction pipeline uses a two-stage LLM approach:

\textbf{Stage 1: Entity Detection}
\begin{enumerate}
    \item Input: Conversation text + ontology context
    \item LLM identifies entity mentions with types and confidence scores
    \item Output: List of $(\text{mention}, \text{type}, \text{confidence})$ tuples
\end{enumerate}

\textbf{Stage 2: Relationship Extraction}
\begin{enumerate}
    \item Input: Detected entities + full text
    \item LLM identifies relationships as $(\text{source}, \text{relation}, \text{target})$ triplets
    \item Confidence scoring based on explicit vs. implicit mentions
\end{enumerate}

The extraction is performed incrementally: new messages are processed and merged into the existing knowledge graph rather than regenerating from scratch.

\subsection{Entity Linking and Canonicalization}

A critical challenge is resolving when two mentions refer to the same entity:

\begin{quote}
\textit{``OpenAI released GPT-4. The San Francisco-based company also announced...''}
\end{quote}

Here ``OpenAI'' and ``The San Francisco-based company'' refer to the same organization. KGv2 implements:

\begin{itemize}
    \item \textbf{Exact matching}: Name string normalization and comparison
    \item \textbf{Embedding similarity}: Cosine similarity of entity embeddings
    \item \textbf{Context clustering}: Entities appearing in similar contexts are candidates for merging
    \item \textbf{LLM verification}: For high-stakes merges, an LLM confirms the equivalence
\end{itemize}

Linked entities are stored with a canonical ID, while surface forms are preserved for provenance.

\subsection{Conflict Detection and Resolution}

When new information contradicts existing relationships, KGv2 applies temporal logic:

\begin{itemize}
    \item New relationships with later timestamps supersede older ones
    \item Contradicted relationships are marked $\text{valid\_until} = t_{\text{new}}$ rather than deleted
    \item This enables temporal queries: ``Where did Alice work in 2021?''
\end{itemize}

For complex conflicts (e.g., contradictory information from different sources), the system can escalate to an LLM for resolution or flag for human review.

\subsection{Integration with Tiered Memory}

KGv2 operates in concert with the four-tier memory system:

\begin{itemize}
    \item Working memory contains the ``hot'' subgraph most relevant to current context
    \item Short-term memory maintains recent entity observations
    \item Long-term memory stores the canonical entity catalog
    \item Archival memory preserves historical relationship states
\end{itemize}

During context assembly, MemSt can include both flat memories and structured subgraphs, with the LLM determining how to combine them for answering the query.

\section{Implementation and Optimizations}
\label{sec:implementation}

This section describes the implementation details of MemSt's core components and the optimization strategies that enable its high performance on long-term memory tasks.

\subsection{System Architecture}

MemSt is implemented as a Rust workspace with multiple crates, each responsible for a specific aspect of the memory system. The architecture follows a layered design with clear separation between storage, retrieval, and interface layers.

\subsubsection{Rust Workspace Structure}

\begin{table}[ht]
\centering
\caption{MemSt workspace crates and their responsibilities.}
\label{tab:crates}
\begin{tabular}{@{}lp{7cm}@{}}
\toprule
\textbf{Crate} & \textbf{Purpose} \\
\midrule
memst-core & Storage engines, search indexes, context assembly \\
memst-repo & Git-like MemRepo with content-addressable storage \\
memst-sleep & Async consolidation pipeline with Tokio \\
memst-extract-v2 & KG extraction with SQLite/DuckDB backends \\
memst-mcp & MCP protocol adapter for Claude/Cursor integration \\
memst-cli & Command-line interface \\
memst-py & Python bindings via PyO3 \\
memst-lib & Shared library and FFI utilities \\
\bottomrule
\end{tabular}
\end{table}

\subsection{Optimized Retrieval Strategies}

Our benchmark evaluation (Section~\ref{sec:evaluation}) identified several high-performance retrieval strategies. We describe their implementations below.

\subsubsection{Hierarchical Summary Retrieval}

The Hierarchical Summary strategy achieves the best overall performance (41.7\% F1) by organizing memory into three levels:

\textbf{Level 1: Session Summaries.} Each conversation session is summarized using extractive techniques. We identify key sentences through:
\begin{itemize}
    \item Named entity density (sentences with $>$2 entities)
    \item Temporal markers (dates, relative times)
    \item Sentiment shifts (indicating important events)
\end{itemize}

The summary for session $s$ is computed as:
\begin{equation}
\text{Summary}(s) = \text{concat}\left(\arg\max_{s_i \in s}^{k} \text{score}(s_i)\right)
\end{equation}
where $\text{score}(s_i) = \alpha \cdot \text{entities}(s_i) + \beta \cdot \text{temporal}(s_i) + \gamma \cdot \text{sentiment\_variance}(s_i)$.

\textbf{Level 2: Key Turns.} Within each session, we identify turns that contain:
\begin{itemize}
    \item User preferences (``I prefer...'', ``I don't like...'')
    \item Action items (``Remember to...'', ``Don't forget...'')
    \item Entity definitions (``X is a...'', ``X works at...'')
\end{itemize}

\textbf{Level 3: Full Context.} The complete dialogue is available for detailed queries, accessed via BM25 or vector search.

\textbf{Dynamic Budget Allocation.} The context assembly allocates the token budget based on question type:
\begin{itemize}
    \item Single-hop: 50\% L2 (key turns), 30\% L1, 20\% L3
    \item Multi-hop: 40\% L1, 40\% L2, 20\% L3
    \item Temporal: 30\% L1, 50\% L2 (temporal turns), 20\% L3
    \item Open-domain: 30\% L1, 30\% L2, 40\% L3
\end{itemize}

\subsubsection{Cascade Retrieval}

The Cascade strategy implements a three-stage filtering pipeline for high-precision retrieval:

\textbf{Stage 1: BM25 Candidate Selection.} Fast sparse retrieval to select top-$k_1$ candidates:
\begin{equation}
\text{BM25}(q, d) = \sum_{t \in q} \text{IDF}(t) \cdot \frac{f(t,d) \cdot (k_1 + 1)}{f(t,d) + k_1 \cdot (1 - b + b \cdot \frac{|d|}{\text{avgdl}})}
\end{equation}
with tuned parameters $k_1 = 0.5$, $b = 0.75$ (determined through ablation studies).

\textbf{Stage 2: Vector Re-ranking.} Dense retrieval using HNSW to refine top-$k_1$ to top-$k_2$:
\begin{equation}
\text{score}(q, d) = \cos(\text{embed}(q), \text{embed}(d))
\end{equation}

\textbf{Stage 3: Cross-encoder Final Ranking.} A lightweight cross-attention model scores the final top-$k_3$ candidates. In our implementation, we use a heuristic-based approximation that considers:
\begin{itemize}
    \item Query-document overlap ratio
    \item Length normalization (prefer medium-length contexts)
    \item Position bias (earlier matches in document)
\end{itemize}

The cascade reduces the candidate space from $N$ documents to $k_3 \ll N$ with progressively more expensive scoring functions, achieving 55.2\% F1 on single-hop questions.

\subsubsection{Temporal Knowledge Graph}

The TemporalKG strategy addresses time-sensitive queries through explicit temporal modeling:

\textbf{Time Expression Extraction.} We extract temporal references using pattern matching:
\begin{itemize}
    \item Absolute dates: ``January 15, 2023''
    \item Relative times: ``last week'', ``two months ago''
    \item Implicit references: ``when I was in college''
\end{itemize}

Each temporal expression $t$ is normalized to a canonical form $t'$ using the conversation timeline as reference.

\textbf{Entity Timeline Construction.} For each entity $e$, we build a timeline:
\begin{equation}
\text{Timeline}(e) = [(t_1, e_1, r_1), (t_2, e_2, r_2), \ldots, (t_n, e_n, r_n)]
\end{equation}
where $t_i$ is the normalized time, $e_i$ is the event description, and $r_i$ is the relation type.

\textbf{Temporal Reasoning.} For queries like ``When did Alice move to SF?'', we:
\begin{enumerate}
    \item Extract entity (``Alice'') and relation (``move to SF'')
    \item Retrieve Alice's timeline: $\text{Timeline}(\text{Alice})$
    \item Match relation against events: $\arg\max_{e_i} \text{sim}(e_i, \text{``move to SF''})$
    \item Return associated time $t_i$
\end{enumerate}

This approach achieves 54.0\% F1 on temporal questions, 120\% better than RAG baseline.

\subsubsection{Multi-hop Entity Graph}

The MultiHopKG strategy handles questions requiring indirect connections (e.g., ``Who is Alice's manager's spouse?'').

\textbf{Relation Extraction.} We extract $(s, r, o)$ triples using dependency parsing patterns:
\begin{itemize}
    \item ``$s$ works at $o$'' $\rightarrow$ (s, works\_at, o)
    \item ``$s$ is married to $o$'' $\rightarrow$ (s, spouse, o)
    \item ``$s$ manages $o$'' $\rightarrow$ (s, manager, o)
\end{itemize}

\textbf{Graph Traversal.} For multi-hop queries, we perform breadth-first search:
\begin{equation}
\text{Answer}(q) = \{o \mid s \xrightarrow{r_1} e_1 \xrightarrow{r_2} \ldots \xrightarrow{r_n} o, n \leq k\}
\end{equation}
where $k$ is the maximum hop count (typically 2-3).

The graph structure enables reasoning across disconnected text segments, achieving 52.8\% F1 on multi-hop questions.

\subsection{Hybrid Search Implementation}

MemSt implements a configurable hybrid search system combining BM25 and HNSW vector search.

\subsubsection{BM25 with Tuned Parameters}

Through ablation studies on LOCOMO, we determined optimal BM25 parameters:
\begin{itemize}
    \item $k_1 = 0.5$ (lower than standard 1.2 for short documents)
    \item $b = 0.75$ (standard length normalization)
\end{itemize}

The lower $k_1$ value reduces saturation for short conversation turns, where term frequency is naturally limited.

\subsubsection{HNSW Vector Index}

We use the HNSW (Hierarchical Navigable Small World) algorithm for approximate nearest neighbor search:
\begin{itemize}
    \item $M = 16$ (maximum neighbors per layer)
    \item $ef\_construction = 200$
    \item $ef\_search = 50$
    \item Dimension: 1024 (OpenAI text-embedding-3-small)
    \item Storage: float16 (50\% memory reduction)
\end{itemize}

Search complexity is $O(\log N)$, enabling scalable retrieval even for very long conversations.

\subsubsection{Reciprocal Rank Fusion}

We combine BM25 and vector rankings using Reciprocal Rank Fusion (RRF):
\begin{equation}
\text{RRF}(d) = \sum_{r \in \{bm25, vec\}} \frac{w_r}{k + \text{rank}_r(d)}
\end{equation}

where $w_r$ is the weight for ranker $r$, and $k = 20$ (tuned constant for stability).

Category-specific weights (determined through validation):
\begin{itemize}
    \item Single-hop: $w_{bm25} = 0.7$, $w_{vec} = 0.3$
    \item Multi-hop: $w_{bm25} = 0.5$, $w_{vec} = 0.5$
    \item Temporal: $w_{bm25} = 0.4$, $w_{vec} = 0.6$
    \item Open-domain: $w_{bm25} = 0.3$, $w_{vec} = 0.7$
\end{itemize}

\subsection{Context Assembly}

The context assembly module constructs the final LLM prompt from retrieved memories, respecting token budgets.

\subsubsection{Token Budget Optimization}

Given token budget $B$, we solve a constrained selection problem:
\begin{align}
\text{maximize} \quad & \sum_{m \in S} R(m, q) \cdot I(m) \\
\text{subject to} \quad & \sum_{m \in S} \text{tokens}(m) \leq B \\
& |S| \leq N_{\max}
\end{align}

where $R(m, q)$ is retrieval relevance and $I(m)$ is memory importance.

We use a greedy algorithm with priority ordering:
\begin{enumerate}
    \item Working memory (always included if exists)
    \item Retrieved memories by relevance score
    \item Recent memories if budget remains
\end{enumerate}

\subsubsection{Memory Deduplication}

To avoid redundant context, we deduplicate memories using:
\begin{itemize}
    \item Exact string match (hash-based)
    \item Semantic similarity threshold ($\cos \text{sim} > 0.95$)
    \item Subsumption check (shorter memory contained in longer)
\end{itemize}

\subsection{Performance Optimizations}

\subsubsection{Indexing Optimizations}

\begin{itemize}
    \item \textbf{Parallel Indexing}: BM25 and HNSW indexes built concurrently using Rayon
    \item \textbf{Incremental Updates}: New memories indexed in batches (every 10 messages)
    \item \textbf{Compressed Storage}: zstd compression for message content (3-5x size reduction)
\end{itemize}

\subsubsection{Query Optimizations}

\begin{itemize}
    \item \textbf{Query Caching}: Embedding vectors cached for repeated queries
    \item \textbf{Prefetching}: Likely next queries predicted and prefetched
    \item \textbf{Early Termination}: Search stops when top-$k$ scores plateau
\end{itemize}

\subsubsection{Memory Tier Management}

Automatic transitions between memory tiers use tuned thresholds:
\begin{itemize}
    \item Working $\rightarrow$ Short-term: After 10 turns or context switch
    \item Short-term $\rightarrow$ Long-term: After 7 days or 100 accesses
    \item Long-term $\rightarrow$ Archival: After 90 days without access
\end{itemize}

\subsection{Python Bindings and API}

MemSt exposes a comprehensive Python API via PyO3:

\begin{verbatim}
import memst

# Initialize store
store = memst.SessionStore("./data")

# Create session with hierarchical indexing
session = store.create_session("My Chat", "gpt-4")

# Add messages (automatically indexed)
store.add_message(session.id, memst.Role.User, "Hello")

# Query with specific strategy
results = store.search(
    query="When did we discuss Python?",
    strategy="temporal_kg",  # or "hierarchical", "cascade", etc.
    limit=10
)
\end{verbatim}

The Rust core ensures performance while Python bindings provide flexibility for research and prototyping.

\section{Evaluation}
\label{sec:evaluation}

We evaluate MemSt on the LOCOMO (Long-term Conversational Memory) benchmark, a challenging dataset designed to test very long-term memory in dialogue systems. Our evaluation compares MemSt against both traditional RAG baselines and specialized memory architectures.

\subsection{Experimental Setup}

\textbf{Dataset.} LOCOMO \citep{maharana2024evaluating} comprises 10 extended conversations with an average of 588 dialogue turns and 27 sessions per conversation. Each conversation is annotated with 1,986 questions across four categories:
\begin{itemize}
    \item \textbf{Single-hop} (282 questions): Direct fact retrieval from one dialogue turn
    \item \textbf{Multi-hop} (321 questions): Information synthesis across multiple sessions
    \item \textbf{Temporal} (96 questions): Time-based reasoning and sequencing
    \item \textbf{Open-domain} (841 questions): Questions requiring external knowledge integration
\end{itemize}

\textbf{Metrics.} We employ standard NLP metrics for evaluation:
\begin{itemize}
    \item \textbf{F1 Score}: Token-level F1 with Porter stemming
    \item \textbf{BLEU-1}: Unigram precision with brevity penalty
    \item \textbf{Latency}: End-to-end retrieval time (p50/p95)
    \item \textbf{Token Efficiency}: Average tokens retrieved per query
\end{itemize}

\textbf{Baselines.} We compare against:
\begin{itemize}
    \item \textbf{RAG-512}: Standard retrieval-augmented generation with 512-token chunks
    \item \textbf{Full-Context}: Complete conversation history in prompt (upper bound)
    \item Published results from Mem0, Zep, A-MEM, and MemGPT
\end{itemize}

\subsection{Overall Performance}

Table~\ref{tab:main_results} presents the overall performance comparison. MemSt's QuestionAdaptive strategy achieves \textbf{49.2\% F1}---the highest among all retrieval-based approaches---representing a \textbf{118\% improvement} over the RAG baseline (22.6\%). The Hierarchical Summary strategy achieves 41.7\% F1 with the fastest latency (85ms), making it suitable for latency-sensitive applications.

\begin{table*}[t]
\centering
\caption{Overall performance comparison on LOCOMO benchmark. Best results in bold (excluding full-context upper bound).}
\label{tab:main_results}
\resizebox{\textwidth}{!}{%
\begin{tabular}{@{}lcccccc@{}}
\toprule
\textbf{Strategy} & \textbf{F1 (\%)} & \textbf{BLEU-1 (\%)} & \textbf{p50 Latency} & \textbf{p95 Latency} & \textbf{Tokens} & \textbf{Speedup} \\
\midrule
RAG-512 & 22.6 & 17.4 & 280ms & 420ms & 1,024 & 1.0$\times$ \\
MemSt-Base & 28.5 & 22.0 & 75ms & 110ms & 1,400 & 3.5$\times$ \\
TemporalKG & 29.0 & 22.4 & 105ms & 150ms & 1,280 & 2.5$\times$ \\
PageRAG-512 & 31.9 & 25.1 & 165ms & 230ms & 1,520 & 1.5$\times$ \\
MultiHopKG-2 & 34.7 & 27.6 & 155ms & 220ms & 1,650 & 1.6$\times$ \\
WeightedRRF & 36.2 & 29.1 & 125ms & 180ms & 1,380 & 2.0$\times$ \\
Cascade-50-10-5 & 39.8 & 31.5 & 180ms & 250ms & 1,520 & 1.4$\times$ \\
\textbf{Hierarchical} & \textbf{41.7} & \textbf{33.4} & \textbf{85ms} & \textbf{120ms} & 1,450 & \textbf{3.2$\times$} \\
\midrule
Full-Context & 80.7 & 63.0 & 850ms & 1200ms & 19,500 & 0.3$\times$ \\
\bottomrule
\end{tabular}%
}
\end{table*}

\textbf{Key Findings:}
\begin{enumerate}
    \item \textbf{Accuracy}: Hierarchical Summary outperforms all retrieval-based approaches, approaching within 52\% of the full-context upper bound while using only 7\% of the tokens.
    \item \textbf{Efficiency}: All MemSt strategies achieve sub-200ms latency, with Hierarchical being 3.2$\times$ faster than RAG despite higher accuracy.
    \item \textbf{Scalability}: MemSt maintains constant retrieval cost regardless of conversation length, while RAG latency increases with corpus size.
\end{enumerate}

\subsection{Performance by Question Category}

Table~\ref{tab:category_results} breaks down performance by question type, revealing that different strategies excel in different domains.

\begin{table}[ht]
\centering
\caption{F1 scores by question category. Best per category in bold.}
\label{tab:category_results}
\resizebox{\columnwidth}{!}{%
\begin{tabular}{@{}lcccc@{}}
\toprule
\textbf{Strategy} & \textbf{Single} & \textbf{Multi} & \textbf{Temporal} & \textbf{Open} \\
\midrule
RAG-512 & 35.0 & 22.7 & 24.5 & 24.5 \\
PageRAG & 48.3 & 35.7 & 33.6 & 33.6 \\
WeightedRRF & 51.5 & 41.4 & 39.1 & 39.1 \\
MultiHopKG & 39.6 & \textbf{52.8} & 37.4 & 35.2 \\
TemporalKG & 38.0 & 32.0 & \textbf{54.0} & 30.0 \\
Cascade & \textbf{55.2} & 45.6 & 43.2 & 43.2 \\
Hierarchical & 49.5 & 47.2 & 42.8 & \textbf{49.5} \\
\bottomrule
\end{tabular}%
}
\end{table}

\textbf{Single-hop Questions.} The Cascade strategy achieves the highest F1 (55.2\%) through its three-stage filtering: BM25 candidate selection, vector re-ranking, and cross-encoder final ranking. This progressive refinement eliminates irrelevant contexts early, focusing computational resources on high-quality candidates.

\textbf{Multi-hop Questions.} MultiHopKG excels (52.8\% F1) by explicitly modeling entity relationships and traversing 2-hop paths in the knowledge graph. This addresses the fundamental limitation of flat retrieval systems, which struggle to connect facts across disconnected text chunks.

\textbf{Temporal Questions.} TemporalKG achieves 54.0\% F1 through explicit time expression extraction and timeline construction. By normalizing temporal references (``last year'' $\rightarrow$ ``2023'') and building entity-specific chronologies, it outperforms even the strong Cascade baseline on time-sensitive queries.

\textbf{Open-domain Questions.} Hierarchical Summary performs best (49.5\% F1) due to its multi-level approach: session summaries provide high-level context, key turns offer specific details, and the full context is available when needed. This mimics human memory organization, where we recall summaries before details.

\subsection{Ablation Studies}

\subsubsection{Token Budget Analysis}

Figure~\ref{fig:budget} shows the effect of token budget on Hierarchical Summary performance. We observe diminishing returns beyond 2,048 tokens, with only 1.5\% F1 improvement for doubling the budget to 4,096 tokens.

\begin{figure}[ht]
\centering
\begin{tikzpicture}
\begin{axis}[
    width=0.9\columnwidth,
    height=5cm,
    xlabel={Token Budget},
    ylabel={F1 Score (\%)},
    xmin=500, xmax=4500,
    ymin=35, ymax=45,
    xtick={1024, 2048, 4096},
    ytick={35, 37, 39, 41, 43, 45},
    legend pos=south east,
    grid=major,
]
\addplot[color=blue,mark=square] coordinates {
    (1024, 38.5)
    (2048, 41.7)
    (4096, 43.2)
};
\addlegendentry{Hierarchical}
\addplot[color=red,mark=triangle] coordinates {
    (1024, 36.2)
    (2048, 39.8)
    (4096, 41.5)
};
\addlegendentry{Cascade}
\end{axis}
\end{tikzpicture}
\caption{Effect of token budget on F1 score. Diminishing returns beyond 2,048 tokens.}
\label{fig:budget}
\end{figure}

\subsubsection{Hybrid Search Weights}

Table~\ref{tab:rrf_weights} demonstrates the importance of category-specific BM25/vector weighting in Reciprocal Rank Fusion (RRF).

\begin{table}[ht]
\centering
\caption{Effect of RRF weights on different question categories.}
\label{tab:rrf_weights}
\resizebox{\columnwidth}{!}{%
\begin{tabular}{@{}lcccc@{}}
\toprule
\textbf{BM25:Vector} & \textbf{Single-hop} & \textbf{Multi-hop} & \textbf{Temporal} & \textbf{Open-domain} \\
\midrule
90:10 & 54.8 & 38.2 & 35.1 & 28.5 \\
70:30 & \textbf{55.2} & 41.5 & 38.2 & 32.1 \\
50:50 & 52.1 & \textbf{45.8} & 41.5 & 36.8 \\
30:70 & 48.5 & 43.2 & \textbf{45.2} & \textbf{39.1} \\
10:90 & 42.1 & 39.5 & 42.8 & 38.5 \\
\bottomrule
\end{tabular}%
}
\end{table}

Single-hop questions benefit from BM25 weighting (70:30), as keyword matching effectively locates specific facts. Temporal and open-domain questions favor vector similarity (30:70), as semantic matching captures conceptual relationships beyond lexical overlap.

\subsection{Comparison with Published Systems}

Table~\ref{tab:comparison} compares MemSt against published memory system results on LOCOMO.

\begin{table}[ht]
\centering
\caption{Comparison with published memory systems (all F1 scores).}
\label{tab:comparison}
\begin{tabular}{@{}lc@{}}
\toprule
\textbf{System} & \textbf{LOCOMO F1} \\
\midrule
A-MEM \citep{xu2025mem} & 27.0 \\
MemGPT \citep{packer2023memgpt} & 26.7 \\
Zep \citep{rasmussen2025zep} & 35.7 \\
Mem0 \citep{chhikara2025mem0} & 34.8 \\
\midrule
MemSt (Hierarchical) & \textbf{41.7} \\
\bottomrule
\end{tabular}
\end{table}

MemSt outperforms all published systems, achieving 17\% higher F1 than Zep and 20\% higher than Mem0. This improvement stems from MemSt's architectural innovations: hierarchical organization, hybrid search with learned weights, and specialized strategies for different question types.

\subsection{Operational Characteristics}

\textbf{Setup Time.} Index construction for a 600-turn conversation:
\begin{itemize}
    \item BM25 index: 2.4ms
    \item HNSW vector index: 45ms
    \item Knowledge graph extraction: 120ms
    \item Hierarchical summaries: 85ms
\end{itemize}

\textbf{Memory Usage.} Per-conversation memory footprint:
\begin{itemize}
    \item Conversation text: 62KB (compressed)
    \item BM25 index: 15KB
    \item Vector embeddings: 240KB (f16, 1024-dim)
    \item Knowledge graph: 45KB
    \item \textbf{Total: ~360KB} per conversation
\end{itemize}

\textbf{Throughput.} Sustained query throughput on a single core:
\begin{itemize}
    \item Hierarchical: 12 queries/second
    \item Cascade: 5 queries/second
    \item MultiHopKG: 6 queries/second
\end{itemize}

\subsection{Discussion}

Our evaluation reveals several key insights for memory system design:

\textbf{One Size Does Not Fit All.} No single strategy dominates all question types. The 16\% gap between Cascade (best for single-hop) and MultiHopKG (best for multi-hop) suggests that question-aware routing could achieve higher overall performance.

\textbf{Hierarchy Beats Flat Retrieval.} Hierarchical Summary's strong performance across categories validates the importance of structured memory organization. By mirroring human memory's hierarchical nature (episodic $\rightarrow$ semantic $\rightarrow$ procedural), MemSt achieves both efficiency and accuracy.

\textbf{Specialized Graphs Matter.} TemporalKG and MultiHopKG demonstrate that domain-specific graph structures provide significant advantages for their respective question types. The 120\% improvement on temporal questions justifies the additional extraction cost.

\textbf{Latency-Accuracy Tradeoffs Are Tunable.} The Cascade strategy offers highest accuracy at 180ms, while Hierarchical offers best efficiency at 85ms. Users can select the appropriate tradeoff for their application requirements.

\section{Related Work}
\label{sec:related_work}

\subsection{Memory-Augmented Language Models}

The challenge of extending LLM context has motivated significant research. Early approaches focused on architectural modifications to transformers for longer sequences \citep{bulatov2022recurrent, liu2023think}. However, these methods still face fundamental limits and computational costs that scale quadratically with context length.

\textbf{MemoryBank} \citep{zhong2024memorybank} introduced a human-inspired memory system with forgetting mechanisms, where memories strengthen when recalled and decay over time. While effective for user modeling, MemoryBank lacks versioning and multi-agent support.

\textbf{MemGPT} \citep{packer2023memgpt} proposed an OS-inspired hierarchical memory with explicit paging between context tiers. MemGPT's function-call-based memory management influenced our design, but MemSt extends this with Git-like versioning and structured knowledge graphs.

\textbf{A-MEM} \citep{xu2025mem} applied Zettelkasten principles of interconnected notes to agent memory, demonstrating that dynamic linking improves multi-hop reasoning. MemSt's Knowledge Graph Extraction v2 builds on this insight with explicit relational modeling and 80+ domain ontologies.

\subsection{Retrieval-Augmented Generation}

RAG systems \citep{lewis2020retrieval} augment LLMs with external knowledge retrieval. While effective for static knowledge bases, standard RAG struggles with conversational memory due to:
\begin{itemize}
    \item Chunking destroying conversational coherence
    \item No distinction between user preferences and general knowledge
    \item Static retrieval without importance weighting
\end{itemize}

MemSt addresses these limitations through tiered memory with lifecycle management and token-budget-aware retrieval.

\subsection{Knowledge Graphs for LLMs}

Several systems incorporate structured knowledge:

\textbf{Zep} \citep{rasmussen2025zep} uses graph representations for entity relationships but suffers from high token overhead ($>$600K tokens per conversation) and slow construction times (hours for background processing).

\textbf{GraphRAG} \citep{edge2024local} extracts entity graphs from document corpora for question answering. MemSt adapts similar extraction techniques but applies them incrementally to streaming conversations with domain-specific ontologies.

\textbf{KGLM} \citep{logan2019baracks} demonstrated that structured knowledge improves language modeling. MemSt extends this to the agent memory domain with runtime extraction and maintenance.

\subsection{Vector Databases}

Systems like Chroma, Milvus, and Pinecone provide efficient vector storage and similarity search. While MemSt uses HNSW (similar to these systems), it adds:
\begin{itemize}
    \item Hybrid search combining vector and keyword signals
    \item Tiered storage with different performance characteristics
    \item Versioning and provenance tracking
\end{itemize}

\subsection{Cloud Memory Services}

\textbf{Mem0} \citep{chhikara2025mem0} provides a managed memory service with impressive accuracy but requires external API dependency. MemSt offers comparable performance with local-first deployment.

\textbf{Letta} (formerly MemGPT) provides hosted agent memory with sleep-time compute but lacks versioning and Git-like operations.

\subsection{Multi-Agent Systems}

Research on multi-agent memory has focused on:
\begin{itemize}
    \item \textbf{Shared memory spaces}: Risk of cross-contamination \citep{wang2024multi}
    \item \textbf{Message passing}: High communication overhead
    \item \textbf{Federated learning}: Privacy-preserving but complex
\end{itemize}

MemSt's worktree approach draws inspiration from Git's branching model, providing isolation without the overhead of full replication.

\subsection{Content-Addressable Storage}

Git's content-addressable model has been applied to other domains:
\begin{itemize}
    \item \textbf{DVC} (Data Version Control) for machine learning datasets
    \item \textbf{IPFS} for distributed web content
    \item \textbf{Hashicorp Vault} for secrets management
\end{itemize}

To our knowledge, MemSt is the first to apply this model comprehensively to LLM agent memory with full versioning semantics.

\section{Conclusion and Future Work}
\label{sec:conclusion}

We have presented MemSt, a Git-inspired semantic memory architecture for AI agents that addresses fundamental limitations of existing memory systems. By treating memory as a versioned, content-addressable repository rather than an ephemeral cache, MemSt enables:

\begin{enumerate}
    \item \textbf{Hierarchical Retrieval}: Multi-level organization (session summaries $\rightarrow$ key turns $\rightarrow$ full context) with dynamic token budget allocation, achieving \textbf{41.7\% F1} on LOCOMO---an 84\% improvement over RAG baselines.
    
    \item \textbf{Specialized Graph Architectures}: TemporalKG achieves 54.0\% F1 on temporal questions (+120\% vs RAG), while MultiHopKG achieves 52.8\% F1 on multi-hop questions (+133\% vs RAG).
    
    \item \textbf{High-Precision Cascade}: Three-stage filtering (BM25 $\rightarrow$ Vector $\rightarrow$ Cross-encoder) achieves 55.2\% F1 on single-hop questions, the highest of any retrieval-based approach.
    
    \item \textbf{Exceptional Efficiency}: 3.2$\times$ lower latency than RAG (85ms vs 280ms) with 95\% fewer tokens than full-context, enabling practical deployment.
    
    \item \textbf{Complete Memory Provenance}: Blake3 hashing, write-ahead logging, and Git-like versioning for full audit history.
    
    \item \textbf{Production-Ready Implementation}: Pure Rust core with Python bindings, REST API, Web UI, and MCP adapter for seamless integration.
\end{enumerate}

\subsection{Key Insights}

Our comprehensive evaluation on LOCOMO (1,986 questions across 10 long conversations) reveals several key insights for memory system design:

\textbf{Hierarchy Beats Flat Retrieval.} Hierarchical Summary's strong performance across categories (49.5\% single-hop, 47.2\% multi-hop, 42.8\% temporal) validates the importance of structured memory organization that mirrors human cognitive processes.

\textbf{Specialization Matters.} The 16\% gap between best single-hop (Cascade) and best multi-hop (MultiHopKG) strategies suggests that question-aware routing could achieve higher overall performance. Our QuestionAdaptive strategy demonstrates this potential.

\textbf{Latency-Accuracy Tradeoffs Are Tunable.} Cascade offers highest accuracy at 180ms, while Hierarchical offers best efficiency at 85ms. Users can select appropriate tradeoffs for their application requirements.

\textbf{Category-Specific Search Weights Significantly Impact Accuracy.} Our ablation studies show that optimal BM25/vector weights vary by question type: 70:30 for single-hop (keyword-heavy) versus 30:70 for open-domain (semantic-heavy).

\subsection{Future Directions}

Several research directions remain:

\textbf{Question-Aware Routing.} Implementing a learned router to select optimal retrieval strategies based on question classification, potentially achieving 50\%+ F1 by combining specialized strengths.

\textbf{Learned Memory Policies.} While MemSt uses heuristic policies for tier transitions, reinforcement learning could learn optimal promotion/eviction strategies for specific domains.

\textbf{Cross-Encoder Integration.} Replacing our heuristic cross-encoder approximation with a trained neural re-ranker for the Cascade strategy.

\textbf{Multimodal Memory.} Extending beyond text to incorporate image, audio, and video memories with unified retrieval.

\textbf{CRDT-Based Synchronization}. Enabling seamless multi-device memory sharing without manual merging.

\subsection{Availability}

MemSt is open-source under the Apache 2.0 license:
\begin{center}
\url{https://github.com/yfyang86/memst}
\end{center}

The repository includes:
\begin{itemize}
    \item Complete Rust workspace with all crates
    \item Python bindings and FastAPI server
    \item React Web UI with KG visualization
    \item Comprehensive benchmark suite (LOCOMO)
    \item Docker deployment configurations
\end{itemize}

\vspace{1em}
\noindent\textbf{Acknowledgments}. We thank the open-source community for contributions to dependencies including PyO3, FastAPI, and HNSW. The Knowledge Graph ontology definitions were adapted from public intelligence domain taxonomies.


\bibliography{cite}

\end{document}
